\section{Conclusion}
\label{sec:conclusion}

We designed and evaluated three CNN-based deepfake detectors on
StyleGAN2-generated face images.
Key findings:

\begin{enumerate}
  \item \textbf{FFT alone is insufficient.}  A ResNet18 on log-magnitude
    spectra achieves only 61\% accuracy (AUC\,=\,0.641) — barely above
    random — suggesting that StyleGAN2 spectral artefacts are not easily
    exploitable by a standard CNN at the scale studied.

  \item \textbf{RGB is a strong baseline.}  Fine-tuning a pretrained
    ResNet18 on raw face images yields 99\% accuracy (AUC\,=\,0.9975) with
    rapid convergence, confirming that ImageNet features transfer well to
    deepfake detection.

  \item \textbf{Hybrid outperforms both modalities.}  Late fusion achieves
    99.5\% accuracy, AUC\,=\,0.9998, and \emph{perfect recall} — the
    spectral branch resolves the residual ambiguous cases left by the
    spatial model.

  \item \textbf{Grad-CAM confirms interpretable attention.}  The RGB branch
    focuses on face-internal regions associated with GAN synthesis errors,
    providing evidence of genuine semantic learning.
\end{enumerate}

\paragraph{Limitations.}
The evaluation is limited to StyleGAN2 faces under clean conditions.
Frequency-based detectors are known to degrade under JPEG
compression~\cite{frank2020frequency}, and cross-generator
generalisation (to Stable Diffusion, FaceSwap, \etc) is untested.

\paragraph{Future work.}
Promising directions include: cross-generator robustness evaluation;
cross-attention fusion to replace naive concatenation; temporal consistency
modelling for video deepfake detection; and quantisation of the Hybrid model
for real-time deployment.

\paragraph{Acknowledgements.}
Dataset by Xhlulu~\cite{dataset140k} (Kaggle).
ImageNet weights via PyTorch \texttt{torchvision}.
Experiments run on a single NVIDIA GPU (CUDA 11.8 / 12.1).
